% =====================================================================
%  VaultProof --- DoraHacks submission (Flare Summer Signal, Bounty 2)
%
%  Compile with:   pdflatex vaultproof-submission.tex
%                  pdflatex vaultproof-submission.tex     (twice, for the TOC)
%
%  Everything you are likely to change lives in one of two places:
%    1. The \newcommand definitions in the FACTS block just below --- the
%       addresses, tx hashes and test counts appear once each and are
%       referenced everywhere else, so updating a number here updates it
%       throughout the document.
%    2. The section bodies, which map one-to-one onto the DoraHacks form
%       fields and are ordered the way the form asks for them.
% =====================================================================

\documentclass[11pt,a4paper]{article}

% ---- encoding & fonts -----------------------------------------------
\usepackage[T1]{fontenc}
\usepackage[utf8]{inputenc}
\usepackage{lmodern}
\usepackage{microtype}

% ---- layout ---------------------------------------------------------
\usepackage[margin=2.4cm]{geometry}
\usepackage{parskip}          % blank line between paragraphs, no indent
\usepackage{titlesec}
\usepackage{enumitem}
\usepackage{fancyhdr}

% ---- tables ---------------------------------------------------------
\usepackage{tabularx}
\usepackage{booktabs}
\usepackage{array}

% ---- misc -----------------------------------------------------------
\usepackage{xcolor}
\usepackage{seqsplit}         % lets long hex strings wrap inside a cell
\usepackage[htt]{hyphenat}    % allow hyphenation inside \texttt identifiers
\usepackage[hidelinks]{hyperref}

% ---- colours --------------------------------------------------------
\definecolor{flare}{HTML}{E62058}
\definecolor{inkmuted}{HTML}{555560}
\definecolor{rulegrey}{HTML}{CCCCD2}

% ---- section styling ------------------------------------------------
\titleformat{\section}
  {\normalfont\Large\bfseries\color{flare}}{\thesection.}{0.6em}{}
\titleformat{\subsection}
  {\normalfont\large\bfseries}{\thesubsection}{0.6em}{}
\titlespacing*{\section}{0pt}{1.6em}{0.6em}
\titlespacing*{\subsection}{0pt}{1.2em}{0.4em}

% ---- running header -------------------------------------------------
\pagestyle{fancy}
\fancyhf{}
\renewcommand{\headrulewidth}{0.4pt}
\fancyhead[L]{\small\color{inkmuted}VaultProof}
\fancyhead[R]{\small\color{inkmuted}Flare Summer Signal --- Bounty 2}
\fancyfoot[C]{\small\thepage}

% Long \code{} identifiers can resist breaking; a little slack lets TeX
% find an acceptable break rather than running into the margin.
\setlength{\emergencystretch}{2em}

% Ragged-right variant of tabularx's X column: justified text in a narrow
% cell produces ugly stretched spacing.
\newcolumntype{L}{>{\raggedright\arraybackslash}X}

% ---- helpers --------------------------------------------------------
% Long hex values: monospace, small, and allowed to wrap mid-string.
\newcommand{\hex}[1]{{\ttfamily\footnotesize\seqsplit{#1}}}
% Inline code / identifiers.
\newcommand{\code}[1]{\texttt{\small #1}}
% A short labelled callout for the "one line" style statements.
\newcommand{\keyline}[1]{%
  \begin{center}\itshape #1\end{center}}

% =====================================================================
%  FACTS --- edit these, not the body text
% =====================================================================
\newcommand{\ProjectName}{VaultProof}
\newcommand{\TeamName}{Prashant Thakur}
\newcommand{\RepoURL}{https://github.com/Lakshay-Deol/VaultProof}
\newcommand{\ChainName}{Flare Coston2}
\newcommand{\ChainID}{114}

% Deployed contracts
\newcommand{\AddrRegistry}{0xD653bE4c296E2462D22254953D2Aaa7D4DA1917C}
\newcommand{\AddrPool}{0x3b7c700cd2d812348de61BD13b28e601C661b5Da}
\newcommand{\AddrSender}{0x45540745B838F6f3feC76E662b5539BcB82339c3}
\newcommand{\AddrMeasurement}{0xe1788fF42Fc5a5B4012d5af6f8B51fe3a3eF36f7}
\newcommand{\AddrUSDC}{0x459c634EE948f6D486b714E06C1F186034F2e7A4}

% Live transactions
\newcommand{\TxAttest}{0xd99378e7a668ff52eef4e032902cd8e2d4df99401b908fb316ac686af73ceb4f}
\newcommand{\TxBorrow}{0x66fde48b7de4c0871b23662fa966e2a13d984d773f90fcb0c0d316e54bd02b46}

% Test counts
\newcommand{\TestsTotal}{73}
\newcommand{\TestsGo}{25}
\newcommand{\TestsFoundry}{11}
\newcommand{\TestsWeb}{33}
\newcommand{\TestsLive}{4}

% =====================================================================

\begin{document}

% ---------------------------------------------------------------- title
\begin{center}
  {\Huge\bfseries \ProjectName\par}
  \vspace{0.4em}
  {\large\color{flare} Prove solvency. Reveal nothing.\par}
  \vspace{0.9em}
  {\color{inkmuted}
   DoraHacks --- Flare Summer Signal\\
   Bounty 2: Confidential Compute Apps\\
   Deployed on \ChainName{} (chain ID \ChainID)\par}
\end{center}

\vspace{0.6em}
{\color{rulegrey}\hrule}
\vspace{1.4em}

\tableofcontents
\vspace{1.2em}

% ================================================================ 1
\section{Project overview}

\begin{tabularx}{\textwidth}{@{}l L@{}}
  \toprule
  \textbf{Project name}   & \ProjectName \\
  \textbf{Bounty}         & Bounty 2 --- Confidential Compute Apps (Flare Confidential Compute) \\
  \textbf{One line}       & Prove solvency. Reveal nothing. \\
  \textbf{Network}        & \ChainName{} testnet, chain ID \ChainID \\
  \textbf{Repository}     & \url{\RepoURL} \\
  \textbf{Team}           & \TeamName \\
  \textbf{Licence}        & MIT \\
  \bottomrule
\end{tabularx}

% ================================================================ 2
\section{Short product description}

\ProjectName{} is a confidential solvency oracle on Flare Confidential Compute.
A borrower seals a read-only exchange API key to a public key that only a
hardware enclave holds. The enclave decrypts it, queries the exchange, prices
the holdings with Flare's own FTSO feeds, and reduces everything to a coarse
solvency tier plus a nullifier and an expiry.

What reaches the chain is a wallet address, a tier, an expiry, a nullifier and
a code-hash measurement --- never the API key, the exchange name, the asset
mix, the exact balance, or the account identifier. A \code{LendingPool}
contract reads that tier to extend undercollateralised credit.

\keyline{FDC brings public data on-chain. FCC brings private conclusions on-chain.}

FDC's security comes from replication, and replication is incompatible with
secrecy --- you cannot ask ten verifiers to check a balance behind one API key
without giving ten verifiers the API key. That gap is what this fills.

% ================================================================ 3
\section{Target user}

\textbf{Crypto-holding borrowers locked out of undercollateralised credit.}
Someone with 2,000 USDC on Flare and \$52,000 on an exchange, who cannot borrow
8,000 USDC because the chain cannot see off-chain wealth and there is no safe
way to prove it.

\textbf{Lending protocols that want a solvency signal without custody of user
data.} A pool that would extend a larger cap given evidence of off-chain
assets, but does not want to hold, process, or be liable for exchange
credentials.

% ================================================================ 4
\section{How it helps day-to-day life}

Maya holds 2,000 USDC on-chain and \$52,000 on Kraken. Her boiler dies; she
needs 8,000 USDC for six weeks. She is good for it many times over --- but
every on-chain lender sees 2,000 USDC and lends her about 1,600.

\begin{tabularx}{\textwidth}{@{}l L@{}}
  \toprule
  \textbf{Her options today} & \textbf{What it costs her} \\
  \midrule
  Sell the exchange position
    & Taxable event, loses a position she wanted to keep, to solve a six-week gap \\
  \addlinespace
  Exchange loan desk
    & KYC, custody, jurisdiction limits, days to weeks \\
  \addlinespace
  Give an API key to a scoring service
    & Their ops team reads her live balances; their breach becomes her breach \\
  \bottomrule
\end{tabularx}

\vspace{0.8em}
Each trades away the position, the time, or the privacy. \ProjectName{} trades
away none: she seals a read-only key, the enclave publishes \code{T3}, the pool
raises her cap.

\subsection{Why it generalises}

The same machine answers other questions. Swap the final reducer and you get:

\begin{itemize}[leftmargin=1.2em, itemsep=0.35em]
  \item \textbf{Proof of income for renting} --- an income band, instead of
        three months of bank statements listing every merchant and habit.
  \item \textbf{Proof of reserves for custodians} --- solvency without
        publishing the treasury layout to competitors and attackers.
  \item \textbf{Private KYC} --- ``over 18, not sanctioned, resident of X'',
        computed over documents no database ever holds.
\end{itemize}

The general shape is \emph{a private fact you can prove but not safely share}.

% ================================================================ 5
\section{Full flow --- six stages}

Each stage is gated on the last in code, not in the UI.

\begin{tabularx}{\textwidth}{@{}c l L L@{}}
  \toprule
  \textbf{\#} & \textbf{Stage} & \textbf{Public} & \textbf{Sealed} \\
  \midrule
  1 & Connect   & wallet address              & --- \\
  \addlinespace
  2 & Verify    & measurement, enclave pubkey & operator cannot substitute a key \\
  \addlinespace
  3 & Seal      & ciphertext length           & API key, secret, exchange \\
  \addlinespace
  4 & Anchor    & keccak256(blob), tx hash    & blob never enters calldata \\
  \addlinespace
  5 & Enclave   & ``processing''              & balance, asset mix, account id \\
  \addlinespace
  6 & Attested  & tier, expiry, nullifier, build & everything else, permanently \\
  \bottomrule
\end{tabularx}

\subsection{The four details that carry the security}

\begin{enumerate}[leftmargin=1.4em, itemsep=0.4em]
  \item \textbf{The signature is verified in the user's browser}, against
        Google's published keys --- not on our backend. A signature our own
        server checks for you proves nothing to you. \code{alg: none} is
        rejected outright.
  \item \textbf{The enclave's public key is inside the signed token}, passed as
        the attestation \code{eat\_nonce}. A relay that swaps in its own key
        fails on the user's machine.
  \item \textbf{Hash first, ciphertext second.} Putting ciphertext in calldata
        publishes it permanently, so a future key compromise would
        retroactively expose every credential ever submitted.
  \item \textbf{Pricing happens inside the enclave, from FTSO.} If it trusted
        the exchange's own USD figure, anyone running a server that speaks the
        Kraken wire format could mint themselves a top tier.
\end{enumerate}

After all six stages the chain knows roughly \textbf{2.5 bits} about Maya's
wealth. That is the entire public footprint of a \$52,000 credit decision.

% ================================================================ 6
\section{How the project uses Flare}

The test of a meaningful integration is what breaks without it, so that is the
column that matters.

\begin{tabularx}{\textwidth}{@{}l L L@{}}
  \toprule
  \textbf{Component} & \textbf{Role} & \textbf{What breaks without it} \\
  \midrule
  FCC / Confidential Space
    & Runs the handler in a VM the host cannot read; the launcher --- outside
      the container's control --- signs a statement of which image booted.
    & There is no version of this on a normal server. Somebody's ops team reads
      live credentials, and the entire product is that nobody can. \\
  \addlinespace
  FTSO
    & Prices holdings inside the enclave, from the chain the loan settles on.
    & A self-hosted ``exchange'' could inflate itself into the top tier. \\
  \addlinespace
  Measurement registry
    & On-chain whitelist of trusted code hashes, checked before encrypting.
    & ``Trust the code, not the operator'' becomes unverifiable. \\
  \addlinespace
  Coston2
    & Hosts registry, pool and sender; anchors requests, enforces caps at
      drawdown.
    & The enclave's output is a dashboard receipt, not a credit decision. \\
  \bottomrule
\end{tabularx}

% ================================================================ 7
\section{What was newly built}

\subsection{Base used as-is}

Flare's \code{fce-extension-scaffold}, \code{flare-foundry-starter} and
\code{flare-viem-starter}.

\subsection{Built during the program}

\begin{itemize}[leftmargin=1.2em, itemsep=0.45em]
  \item \textbf{Five Solidity contracts}, deployed and Blockscout-verified.
        \code{SolvencyRegistry} enforces expiry in the read path, handles
        nullifier collisions, and re-checks the measurement at write time.
        \code{LendingPool} applies tier caps that extend rather than replace
        collateral. \code{VaultProofInstructionSender} handles the anchor, the
        TEE write path and FCC routing. Plus the measurement registry and a
        test USDC token.

  \item \textbf{The solvency enclave} in Go --- HPKE unseal, Kraken
        HMAC-SHA512 signed reads, FTSO pricing resolved through the Flare
        contract registry, HMAC nullifier, zeroize in defers. Credentials are
        held in \code{[]byte}, never \code{string}, because Go strings cannot
        be scrubbed.

  \item \textbf{Real hardware attestation} --- the enclave requests a token
        from the Confidential Space launcher over its unix socket, binds its
        X25519 key as the \code{eat\_nonce}, and reports the signed
        \code{image\_digest} as its measurement. A fresh token per request,
        not one minted at boot and replayed.

  \item \textbf{Real quote verification} --- the browser fetches Google's
        Confidential Computing key set and verifies the token signature with
        WebCrypto, enforcing issuer and expiry.

  \item \textbf{The whole frontend} --- 4 pages, 6 stages, mock/live adapter
        seams, a state machine, real in-browser HPKE, and the design system.
        The directory was empty at the start of the program.

  \item \textbf{Both live adapters} --- viem/wagmi against the deployed
        contracts, and the enclave client against \code{/quote} and
        \code{/sealed}.

  \item \textbf{The attack simulator} --- stage 2 ships a ``Try breaking it''
        control that runs a substituted relay key and an unwhitelisted build
        against the same verifier the honest path uses. Both hard-stop, with
        no override.

  \item \textbf{\TestsTotal{} tests} --- \TestsGo{} Go, \TestsFoundry{}
        Foundry, \TestsWeb{} web unit, \TestsLive{} live against Coston2.
\end{itemize}

\subsection{Three real bugs found and fixed}

\begin{enumerate}[leftmargin=1.4em, itemsep=0.4em]
  \item The scaffold's \code{teeutils.ToHash} right-pads a string into a
        \code{bytes32}, while the contract and frontend use \code{keccak256}.
        The two disagree completely and the failure is silent --- the exact
        vanishing-instruction bug the spec warns about.
  \item Blob delivery was keyed by SHA-256 while the browser anchors
        \code{keccak256}, so the enclave would never have found the blob.
  \item Attestation check 2 verified the token's \emph{binding} but never its
        \emph{signature}, so a forged token with matching fields would have
        passed all three checks.
\end{enumerate}

% ================================================================ 8
\section{Deployment details}

\ChainName{} testnet, chain ID \ChainID. All five contracts verified on
Blockscout.

\begin{tabularx}{\textwidth}{@{}l L@{}}
  \toprule
  \textbf{Contract} & \textbf{Address} \\
  \midrule
  SolvencyRegistry            & \hex{\AddrRegistry} \\
  LendingPool                 & \hex{\AddrPool} \\
  VaultProofInstructionSender & \hex{\AddrSender} \\
  TeeMeasurementRegistry      & \hex{\AddrMeasurement} \\
  MockUSDC (tUSDC)            & \hex{\AddrUSDC} \\
  \bottomrule
\end{tabularx}

\subsection{Live evidence}

\begin{itemize}[leftmargin=1.2em, itemsep=0.35em]
  \item T3 attestation delivered --- \hex{\TxAttest}
  \item Borrowed 15,000 tUSDC --- \hex{\TxBorrow}
  \item A 50,000 attempt reverted with \code{over cap} --- the contract refuses
        rather than trusting the front end.
\end{itemize}

% ================================================================ 9
\section{Technical execution}

Each part can be checked independently:

\begin{quote}
\ttfamily\small
cd web \&\& npm run dev\\
\hspace*{2em}\textrm{\itshape localhost:3000, mock mode, no wallet or funds needed}\\[0.3em]
cd extension/go \&\& go run ./cmd\\
\hspace*{2em}\textrm{\itshape the enclave on :8080}\\[0.3em]
curl localhost:8080/quote\\
\hspace*{2em}\textrm{\itshape measurement, enclave pubkey, attestation mode}\\[0.3em]
go run ./cmd/smoke\\
\hspace*{2em}\textrm{\itshape seals as the browser does, attests, checks replay}\\[0.3em]
cd web \&\& npm run test:live\\
\hspace*{2em}\textrm{\itshape \TestsLive{} tests against the real Coston2 deployment}
\end{quote}

\code{cmd/smoke} proves the two HPKE implementations agree: it fails at
\code{unseal:} if they ever drift, and at \code{query:} if the credential
simply is not a real Kraken key --- so the message tells you which.

% ================================================================ 10
\section{Honest limitations}

\begin{itemize}[leftmargin=1.2em, itemsep=0.45em]
  \item \textbf{Hardware mode has never executed.} The attestation code is real
        and on the real path --- on Confidential Space it fetches a genuine
        AMD SEV-SNP-backed token and the browser verifies it against Google.
        But it has not been run on that hardware, so no claim is made that it
        has. Locally it reports simulated mode and the UI renders it as such.
        Every failure branch fails closed to simulated rather than
        overclaiming.

  \item \textbf{Attestation freshness.} A quote proves the code that booted,
        not the code running at request time. Per-request tokens narrow this;
        nothing closes it.

  \item \textbf{Side channels.} AMD SEV has a published history of them.
        Adequate at moderate size; not a claim of perfection.

  \item \textbf{Exchange trust.} If Kraken lies about a balance, \ProjectName{}
        faithfully attests a lie. The system proves what a named source said,
        not what is true.

  \item \textbf{Governance.} Whoever controls whitelisting is a trusted role.
        The system is not trustless end to end and does not claim to be.
\end{itemize}

% ================================================================ 11
\section{Roadmap}

\begin{enumerate}[leftmargin=1.4em, itemsep=0.35em]
  \item \textbf{Run it on Confidential Space} --- build reproducibly, list the
        resulting image digest on the measurement registry, deploy, and serve
        hardware mode.
  \item \textbf{Reproducible-build verifier CLI} --- so strangers can rebuild
        the image and check the hash themselves.
  \item \textbf{Multi-TEE quorum} --- use the cosigner threshold already in the
        FCC instruction params; require $k$-of-$n$ enclaves to agree.
  \item \textbf{Revocation} --- an explicit revoke path and a kill switch for a
        superseded build.
  \item \textbf{More sources} --- Coinbase, Binance, then bank APIs behind one
        credential schema.
  \item \textbf{New reducers} --- proof of income, proof of reserves, private
        KYC assertions.
  \item \textbf{Governance hardening} --- whitelist control to a multisig, then
        to a lender-governed set.
  \item \textbf{Audit, Songbird, mainnet.}
\end{enumerate}

% ================================================================ 12
\section{Distribution and testing}

No user acquisition yet --- the project is pre-public-demo. Testing to date is
the \TestsTotal-test suite plus on-chain runs against the live deployment. The
honest statement is that this has been validated technically, not yet with real
users.

% ================================================================ 13
\section{Repository, team and licence}

\begin{tabularx}{\textwidth}{@{}l L@{}}
  \toprule
  Repository    & \url{\RepoURL} \\
  Design spec   & \code{docs/design-spec.pdf} --- threat model \S10, build plan
                  \S11, demo script \S12 \\
  Team          & \TeamName \\
  Licence       & MIT \\
  \bottomrule
\end{tabularx}

\vspace{1.4em}
{\color{rulegrey}\hrule}
\vspace{0.8em}

\noindent
{\small\color{inkmuted}
% ---- TODO before submitting: fill these two in ----------------------
\textbf{Still to fill in:} the public demo URL and the recorded video
walkthrough.}

\end{document}
